\chapter{Klasa std::string}
Do sada smo u našim programima koristili C-ovske stringove, odnosno nizove znakova (charova) koji na kraju imaju \texttt{'\textbackslash 0'} (nul-znak), a za rad s njima koristili smo funkcije iz zaglavlja \texttt{string.h} koje su primale stringove kao argument. 
U slučaju C-ovskih stringova, mora voditi računa o veličini polja, pravilnoj alokaciji memorije te mogućim pogreškama poput prekoračenja granica niza pa vrlo lako dolazi do pogrešaka. 

Kako bi se pojednostavio rad s tekstualnim podacima, C++ uvodi klasu \texttt{std::string} koja predstavlja stringove. Klasa \texttt{std::string} automatski upravlja memorijom te sadrži velik broj metoda za stvaranje, spajanje, uspoređivanje, pretraživanje i mijenjanje stringova. Zbog toga je u većini slučajeva preporučljivo koristiti \texttt{std::string} umjesto C-ovskih stringova.

Objasnimo prvo znaćenje riječi klasa. \dots Klasama ćemo se detaljnije baviti u poglavlju \ref{}.

U nastavku dajemo pregled nekih metoda klase \texttt{std::string}. Popis i opis svih metoda nalazi se \href{https://cplusplus.com/reference/string/string/}{ovdje}.

Za korištenje klase \texttt{std::string} potrebno je uključiti zaglavlje \texttt{<string>}:
\begin{minted}{C++}
    #include <string>
\end{minted}

Pogledajmo prvo najčešće konstruktore.
\begin{minted}{C++}
    std::string S1;
    std::string S2("Racunarski praktikum 1");
    std::string S3(5, 'Z');
\end{minted}
Prva linija stvara prazan string, druga linija stvara string \texttt{"Racunarski praktikum 1"}, dok treća linija stvara string \texttt{"ZZZZZ"} (dakle, 5 string koji se sastoji od pet znakova \texttt{'Z'}).

Ekvivalenstno drugoj liniji je 
\begin{minted}{C++}
    std::string S2;
    S2 = std::string("Racunarski praktikum 1");
\end{minted}
te
\begin{minted}{C++}
    std::string S2 = "Racunarski praktikum 1";
\end{minted}




\begin{zadatak}
    U zadatku~\ref{zad:student-dinamicki} C-ovske stringove zamijenite klasom \texttt{std::string}.
\end{zadatak}

\begin{rjesenje}
    Datoteka \texttt{student.h}
    \begin{minted}[
        linenos,
        numbersep=5pt
    ]{C++}
        #pragma once
        #include <string>

        struct Student{
            std::string ime, prezime, jmbag;
            std::string *kolegiji;
            unsigned *ocjene, n;

            Student();
            Student(std::string _ime, std::string _prezime, std::string _jmbag):

            bool ocijeni(std::string kolegij, unsigned ocjena);
            void info();

            ~Student();
        };
    \end{minted}
    Datoteka \texttt{student.cpp}

    \begin{minted}[
        linenos,
        numbersep=5pt
    ]{C++}
        #include <iostream>
        #include <string>
        #include "student.h"

        using namespace std;
    \end{minted}

    Konstruktori:
    \begin{minted}[
        linenos,
        numbersep=5pt
    ]{C++}
        Student::Student(){
            kolegiji = ocjene = nullptr;
            n = 0;
        }

        Student::Student(string _ime, string _prezime, string _jmbag){
            ime = _ime;
            prezime = _prezime;
            jmbag = _jmbag;
            kolegiji = ocjene = nullptr;
            n = 0;
        }
    \end{minted}

    Funkcija \texttt{ocijeni}

     \begin{minted}[
        linenos,
        numbersep=5pt
    ]{C++}
         bool Student::ocijeni(string kolegij, unsigned ocjena){
            int i;
            for(i = 0; i < n; ++i){
                if(kolegiji[i] == kolegij) break;
            }
            if(i < n){ // kolegij vec postoji
                ocjene[i] = ocjena;
                return false;
            }

            unsigned ocjene_kopija = new unsigned[n+1];
            char *kolegiji_kopija = new char[21][n+1];
            for(int i = 0; i < n; ++i){
                ocjene_kopija[i] = ocjene[i];
                kolegiji_kopija[i] = kolegiji[i];
            }
            ocjene[n] = ocjena;
            novi_kolegiji[n] = kolegij;
            delete[] ocjene;
            delete[] kolegiji;
            ocjene = ocjene_kopija;
            kolegiji = kolegiji_kopija;
            ++n;
            return true;
        }
    \end{minted}
\end{rjesenje}

\begin{zadatak} \label{zad:split-string}
    Napišite glavni program koji učitava podatke o studentima u sljedećem obliku:
    \begin{minted}[
        linenos,
        numbersep=5pt
    ]{C++}
        ime_1,prezime_1,jmbag_1
        ime_2,prezime_2,jmbag_2
        ...
        ime_n,prezime_n,jmbag_n
        kraj
    \end{minted}
    \texttt{ime\textunderscore i,prezime\textunderscore i,jmbag\textunderscore i} je jedan string! Ime, prezime i JMBAG studenta međusobno su odvojeni sa \mintinline{bash}{,} (zarezom).

    Podaci se učitavaju sve dok se ne unese string \texttt{"kraj"} budući da broj studenata nije poznat unaprijed. Možete pretpostaviti da će stringovi biti u zadanom formatu.
\end{zadatak}

\begin{rjesenje}
    \begin{minted}[
        linenos,
        numbersep=5pt
    ]{C++}
        #include <iostream>
        #include "student.h"
        #include "fakultet.h"

        using namespace std;
        
        int main(void){
            Fakultet F;
            string str;
            while(1){
                cout << "Unesite podatke o studentu: ";
                cin >> str;
                if(str == "kraj") break;
                string ime, prezime, jmbag;
                ime = str.substr((str.find(",")));
                str.erase(0, ime.size()+1);
                prezime = str.substr((str.find(",")));
                str.erase(0, prezime.size()+1);
                jmbag = str;
                F.dodajStudenta(ime, prezime, jmbag);
            }
            return 0;
        }
    \end{minted}
\end{rjesenje}

\begin{razmisljanje}
    \begin{enumerate}
        \item Riješite zadatak~\ref{zad:split-string} bez pretpostavke da je uneseni string u ispravnom formatu, odnosno dodajte provjeru ispravnosti.
        \item Riješite zadatak~\ref{zad:split-string} uz pretpostavku da se svi podaci o studentima unose kao jedan string u sljedećem obliku:
         \begin{center}
             \texttt{ime\textunderscore 1,prezime\textunderscore 1,jmbag\textunderscore 1;...;ime\textunderscore n,prezime \textunderscore n,jmbag\textunderscore n}
         \end{center}   
         Dakle, podaci o različitim studentima odvojeni su sa \mintinline{bash}{;} (točka zarez).
         \item Implementirajte strukturu \texttt{MyString} koja ima dio funkcionalnosti klase \texttt{std::string}:
                \begin{itemize}
                    \item sve konstruktore koje smo spomenuli u ovom poglavlju
                    \item funkcije \texttt{size}, \texttt{append}, \texttt{substr}, \texttt{find}
                    \item umjesto \texttt{A = B} neka postoji funkcija \texttt{A.copy(B)}
                    \item umjesto \texttt{A == B} neka postoji funkcija \texttt{A.isEqual(B)}
                    \item umjesto \texttt{A < B} neka postoji funkcija \texttt{A.isLessThan(B)}
                \end{itemize}
    \end{enumerate}
\end{razmisljanje}